% !TEX root = ../b-rg.tex
\chapter{Nouns and the Noun Phrase}\label{ch:nouns}
\stub

% --- Planned sections ---
% \section{The noun}
%   No gender, number or case (history: loss of final -s, note 121.5 and the
%   historical appendix); count vs mass; proper nouns; sour 'sister' as the
%   one surviving old nominative.
% \section{Marking number}
%   Number is marked on determiners and by verb agreement, not on the noun:
%   l'oc scuir pasc / l'ec scuir pascn. Unmarked NPs are number-neutral.
%   Cross-reference Articles and Demonstratives.
% \section{Structure of the noun phrase}
%   (possessive) (article) (demonstrative) (numeral/quantifier) N
%   (adjective) (PP) (relative clause); a possessive excludes the article
%   before a consonant and fuses with it before a vowel (my fraðr, mell'orin;
%   see Articles and Demonstratives); adjectives mostly postnominal;
%   stacked modifiers of the same type; emphatic meðes; adverbs inside
%   the NP.
% \section{Genericity}
%   How generic reference is expressed (bare noun? article?).
% \section{Possession and other adnominal relations}
%   de + NP (dy, dell'); possessive determiners (see Pronouns);
%   noun + noun juxtaposition vs de-phrases.
% \section{Nominalisation}
%   Infinitives and adjectives used as nouns; pointer to Word Formation.
